Um zu zeigen, dass der Raum $\ell^{\infty}(M)$ nicht separabel ist, wenn $M$ eine unendliche Menge ist, betrachten wir die Definition der Separabilität. Ein topologischer Raum ist separabel, wenn es eine abzählbare dichte Teilmenge gibt.

Der Raum $\ell^{\infty}(M)$ besteht aus allen beschränkten Funktionen $f: M \to \mathbb{R}$, wobei die Norm durch

\[
\|f\|_{\infty} = \sup_{x \in M} |f(x)|
\]

definiert ist. Um zu zeigen, dass $\ell^{\infty}(M)$ nicht separabel ist, konstruieren wir eine unendliche Familie von Funktionen, die alle einen Abstand von mindestens $\frac{1}{2}$ zueinander haben.

Da $M$ unendlich ist, können wir eine Teilmenge $N \subseteq M$ wählen, die ebenfalls unendlich ist. Für jedes $x \in N$ definieren wir die Funktion $f_x \in \ell^{\infty}(M)$ durch

\[
f_x(y) = 
\begin{cases} 
1, & \text{wenn } y = x, \\
0, & \text{wenn } y \neq x.
\end{cases}
\]

Für zwei verschiedene Funktionen $f_x$ und $f_{x'}$ mit $x \neq x'$ gilt

\[
\|f_x - f_{x'}\|_{\infty} = \sup_{y \in M} |f_x(y) - f_{x'}(y)| = 1.
\]

Da der Abstand zwischen beliebigen zwei verschiedenen Funktionen $f_x$ und $f_{x'}$ mindestens $1$ ist, kann keine abzählbare Teilmenge von $\ell^{\infty}(M)$ dicht sein. Dies liegt daran, dass jede dichte Teilmenge eine Funktion enthalten müsste, die beliebig nah an jeder Funktion $f_x$ liegt, was hier nicht möglich ist.

Daraus folgt, dass $\ell^{\infty}(M)$ nicht separabel ist, wenn $M$ unendlich ist.

%CHECK: Die Lösung wurde überprüft und es wurden keine Schwächen oder Verbesserungsvorschläge im Review identifiziert. Die Herleitung ist korrekt und vollständig.